% page 171 of the facsimile (image p-170 of the scan)
% block 167.6 x 233.3 mm ; left margin 34.4 mm
\pgc{12.700mm}{0.000mm}{
\l{\cel{54}161}
\l{}
\l{}
\l{\sou{fekulo.}\cel{2}Sedimento amilatra quan depozas la suki di ula materii}
\l{\cel{3}vejetala, o quan formacas la farino lavita di la grani de}
\l{\cel{3}cereali, leguminosi, di la tuberkuli de terpomi (= kartofli),}
\l{\cel{3}di la rizomi de manioko. - EFIS.}
\l{}
\l{\sou{fekunda}.\cel{2}Apta produktar enti simila a su. (\sou{anke metaf.})- EFIS.}
\l{}
\l{\sou{felo}. (\sou{tekn}.)Pelo deprenita de animalo mortinta, e preparita por}
\l{\cel{3}skopi diversa. - DE.}
\l{}
\l{\sou{feldo}. (\sou{tekn.}) Spaco, areo, en qua eventas ula distributo di}
\l{\cel{3}forci (elektrala, magnetala, e c.)\cel{2}- FE.}
\l{}
\l{\sou{felgo}. (\sou{tekn.}) I. Singla ek la sis peci de ligno kurvigita qua,}
\l{\cel{3}en roto, formacas la cirklo-kurvo extera en qua inkastresis}
\l{\cel{3}la spoki. - II. Singla ek la quar peci de ligno, plu granda,}
\l{\cel{3}qui formacas cirklo-kurvo duesma, tenono-asemblita en la}
\l{\cel{3}unesma. - DE.}
\l{}
\l{\sou{felica}. Di qua la anmo, la psiko, esas satisfacita komplete.EFIS .}
\l{}
\l{\sou{felina}. Di qua la movi e gesti esas dolca e flexebla quale olti}
\l{\cel{3}di la kato. - EFIS.}
\l{}
\l{\sou{felpo}. Stofo de lano, de kotono, de silko qua imitas plusho,}
\l{\cel{3}ma di qua la pili esas plu longa e min densa. - DI.}
\l{}
\l{\sou{felto}. Texuro forta, densa, ek lano e pili, aglutinita e fulita.}
\l{\cel{3}- DEFIS.}
\l{}
\l{\sou{femina}. Qua apartenas a la sexuo organizita por konceptar e par-}
\l{\cel{3}turar, o por\cel{2}pondar ovi. - DEFIRS.}
\l{}
\l{\sou{femuro}. (\sou{anat.)}\cel{2}Osto di la kruro. - EFIS.}
\l{}
\l{feno.\cel{2}Herbo quan onu falchas e lasas sikeskar, uzata kom nutrivo}
\l{\cel{3}por la kavali, la brutaro. - FIS.}
\l{\cel{13}2 5\cel{2}6 4\cel{11}3}
\l{\sou{fenacetino}. C H .C H .NH. DC. CH . -\cel{2}DEF.}
\l{}
\l{\sou{fenco.}\cel{2}Barilo qua klozas. - DE.}
\l{}
\l{\sou{fendar}. (\sou{trans.}) I. Dividar longesale. - II. Dividar per eskartar}
\l{\cel{3}la parti poke adheranta. - FIS.}
\l{}
\l{\sou{fenestro}. I. Aperturo di qua la formo e la grandeso esas diversa,}
\l{\cel{3}facita en la muri di konstrukturo por lasar aero e jorno pene-}
\l{\cel{3}trar. - II. Kadro vitroplakoza qua klozas ta aperturo. - DEFI.}
\l{}
\l{\sou{fen}ikulc. (\sou{bot.}) Planto aromata de la familio \textquotedbl{}umbeliferi\textquotedbl{}.}
\l{\cel{3}L.\cel{1}\sou{foeniculum vulgare}. - DEFIL.}
}
